% =============================================================================
%  O modelo de objeto por dentro
%  Cobre os subprojetos layout, type_erasure e linkage.
% =============================================================================

% =============================================================================
\section{O modelo de objeto por dentro}
\label{sec:modelo-objeto}
% =============================================================================

Os capítulos anteriores trataram do que a linguagem \emph{diz}. Este trata do
que ela \emph{faz} --- byte a byte, símbolo a símbolo. São três perguntas, cada
uma com um subprojeto para respondê-la com números em vez de afirmações:

\begin{enumerate}
    \item \textbf{Como os bytes de um objeto são arrumados?} O
          \projeto{layout} responde: alinhamento, \emph{padding}, as categorias
          de tipo que autorizam \texttt{memcpy}, e o que \texttt{virtual}
          acrescenta ao objeto.
    \item \textbf{Como se guarda um objeto sem mencionar o tipo dele?} O
          \projeto{type\_erasure} constrói a \classe{std::function} do zero,
          mede a otimização de \emph{buffer} pequeno com um \texttt{operator
          new} instrumentado e compara quatro formas de despacho no cronômetro.
    \item \textbf{Como o compilador e o \emph{linker} combinam sobre nomes,
          tamanhos e endereços?} O \projeto{linkage} --- único subprojeto do
          repositório com \emph{várias} unidades de tradução --- mede a One
          Definition Rule com dois \texttt{.cpp} reais e chama código C
          compilado por um compilador C de verdade.
\end{enumerate}

A ordem não é arbitrária: o apagamento de tipo só faz sentido depois de se saber
o que um ponteiro de \gls{vtable} custa em bytes, e a \gls{abi} da última seção
é, em boa medida, a soma das duas primeiras --- \emph{layout} de objeto mais
nomes de símbolo.

\begin{porque}[por que o modelo de objeto é assunto de C++, e não de Java]
O C++ é a única linguagem de uso corrente em que o programador \emph{escolhe} o
modelo de objeto e paga por ele. Em Java ou Python todo objeto tem cabeçalho,
toda chamada de método é virtual e todo valor mora no \emph{heap}: não há
decisão a tomar. Em C++, um destrutor virtual que ninguém usa custa 8~bytes por
objeto --- 76~MB num vetor de dez milhões de partículas, medidos na
\cref{sec:layout-vptr} --- e trocar a ordem de quatro membros corta um terço do
tamanho da \texttt{struct}. Quem não conhece o modelo não faz nenhuma dessas
escolhas: só sofre as consequências.
\end{porque}

% =============================================================================
\subsection{O layout de um objeto}
\label{sec:layout}
% =============================================================================

\subsubsection{Alinhamento: a exigência do hardware}

Uma definição primeiro. \textbf{Alinhamento} é a exigência de que o endereço de
um objeto seja múltiplo de um certo número. Todo tipo \texttt{T} tem um
\texttt{alignof(T)}, e o objeto só pode começar em endereço divisível por ele:
um \texttt{double} com \texttt{alignof}~8 mora em $0, 8, 16, \ldots$, nunca
em~3.

A exigência não vem da linguagem, vem do processador. A memória é lida em blocos
de tamanho fixo --- a \textbf{linha de cache}, tipicamente 64~bytes em x86-64
--- e as instruções de carga esperam operandos que não atravessem fronteiras.
Ler um \texttt{double} desalinhado custa duas leituras em vez de uma no x86, e
em ARM antigo aborta o processo com \texttt{SIGBUS}.

\begin{regra}[as três regras do layout]
\textbf{(1)}~Todo tipo tem \texttt{alignof(T)}: o endereço precisa ser múltiplo
dele. \textbf{(2)}~O alinhamento de uma \texttt{struct} é o \textbf{maior} entre
os de seus membros. \textbf{(3)}~\texttt{sizeof} é sempre múltiplo de
\texttt{alignof}, para que \emph{arrays} funcionem. Todo o resto --- inclusive
por que \texttt{sizeof} quase nunca é a soma dos membros --- decorre daí.
\end{regra}

O nome do ``resto'' é \textbf{\emph{padding}}: bytes que o compilador insere
entre os membros, e depois do último, para que cada um caia em endereço válido.
O \emph{padding} não guarda nada, não faz parte do valor do objeto e --- como se
verá --- pode conter qualquer coisa.

\subsubsection{A ordem dos membros muda o tamanho}

\begin{porque}[por que declarar em ordem decrescente de alinhamento]
O compilador \textbf{não pode} reordenar os membros de dados: a ordem de
declaração é garantida pelo padrão (para membros com o mesmo especificador de
acesso), justamente para que o \emph{layout} seja previsível e interoperável com
C. A única pessoa que pode arrumar o \emph{padding} da sua \texttt{struct} é
você.
\end{porque}

\begin{lstlisting}[style=bt, language=C++,
    caption={Os mesmos membros, duas ordens, em
             \texttt{src/layout/demo\_tamanho.hpp}.}]
struct Ruim {          // 1 + 8 + 1 + 4 = 14 bytes de dado...
    char      a;
    double    b;
    char      c;
    int       d;
};

struct Bom {           // os MESMOS membros, decrescendo por alinhamento
    double    b;
    int       d;
    char      a;
    char      c;
};
\end{lstlisting}

\begin{lstlisting}[style=bt, language=bash,
    caption={Saída de \texttt{src/layout/}: 14 bytes de dado útil, 8 de
             diferença.}]
struct Ruim { char a; double b; char c; int d; };
  sizeof  = 24   alignof = 8
  offsetof: a=0  b=8  c=16  d=20
  mapa: [a][×7 pad][b b b b b b b b][c][×3 pad][d d d d]

struct Bom  { double b; int d; char a; char c; };
  sizeof  = 16   alignof = 8
  mapa: [b b b b b b b b][d d d d][a][c][×2 pad]

MESMOS dados, 8 bytes de diferença (33% menor).
Num vetor de 1 milhão: 22.9 MB contra 15.3 MB.
\end{lstlisting}

O \texttt{char~a} de \texttt{Ruim} ocupa o \emph{offset}~0 e o \texttt{double~b}
seguinte só pode começar no~8: sete bytes perdidos. Depois o \texttt{char~c}
no~16 e o \texttt{int~d} no~20: mais três. Em \texttt{Bom} nada disso acontece,
porque os membros grandes vêm primeiro e os pequenos preenchem o rabo. Trinta e
três por cento de memória de graça --- e, o que costuma importar mais, um terço
a menos de linhas de cache tocadas em cada passada.

\subsubsection{Por que existe \emph{padding} no fim}

\texttt{struct \{ int a; char b; \}} soma cinco bytes e ocupa oito. A razão é a
regra~(3), e a razão da regra~(3) são os \emph{arrays}: em \texttt{ComRabo v[3]}
cada elemento precisa começar em múltiplo de \texttt{alignof}. Se
\texttt{sizeof} fosse~5, \texttt{v[1]} cairia no \emph{offset}~5 --- e o
\texttt{int} de dentro dele, desalinhado.

\subsubsection{\texttt{memcmp} e o lixo do \emph{padding}}

O subprojeto constrói duas instâncias com os \textbf{mesmos valores de membro},
uma delas com o \emph{padding} marcado com \texttt{0xAA}:

\begin{lstlisting}[style=bt, language=bash,
    caption={Saída de \texttt{src/layout/}: membros iguais, bytes diferentes.}]
x.a == y.a && x.b == y.b     → true
memcmp(&x, &y, sizeof)  == 0 → false
  x: 07 00 00 00 7A AA AA AA
  y: 07 00 00 00 7A 00 00 00
                  ^^ ^^ ^^   ← padding
\end{lstlisting}

\begin{armadilha}[\texttt{memcmp} compara bytes; \texttt{==} compara valores]
O conteúdo do \emph{padding} é \textbf{indeterminado}: ele não faz parte do
valor do objeto e nada obriga o compilador a inicializá-lo. Daí três
proibições: nunca comparar \texttt{struct}s com \texttt{memcmp} (use
\texttt{operator==}, escrito ou \texttt{= default}); nunca transmitir nem
persistir uma \texttt{struct} com \emph{padding} sem zerá-la antes --- o
\emph{padding} \textbf{vaza memória do processo} para o disco ou para a rede,
que é a família de bug do Heartbleed em outra forma; e, se precisa de bytes,
serializar campo a campo.
\end{armadilha}

\subsubsection{\texttt{struct} vazia: por que 1 e não 0}

\texttt{sizeof(struct\{\})} é \textbf{1}. Zero seria mais honesto e quebraria a
linguagem: dois objetos distintos precisam de endereços distintos, e num
\texttt{Vazia v[3]} os três ocupariam o mesmo lugar. Métodos não ocupam espaço:
uma função membro é uma função comum com um \texttt{this} implícito, guardada
uma vez no binário --- só \texttt{virtual} acrescenta algo ao objeto
(\cref{sec:layout-vptr}). E há a exceção que salva o dia:
\texttt{struct HerdaVazia : Vazia \{ int x; \}} tem \texttt{sizeof}
\textbf{4}, não~8, porque a base vazia foi otimizada para zero bytes --- a EBO
da \cref{sec:layout-ebo}.

\begin{table}[htbp]
\centering
\caption{\texttt{sizeof} e \texttt{alignof} medidos por \texttt{src/layout/}
         (GCC 13, x86-64 System V).}
\label{tab:layout-std-tipos}
\footnotesize
\begin{tabular}{@{}lrr@{}}
\toprule
\textbf{Tipo} & \textbf{\texttt{sizeof}} & \textbf{\texttt{alignof}} \\
\midrule
\texttt{int}                 & 4  & 4  \\
\texttt{double}              & 8  & 8  \\
\texttt{long double}         & \textbf{16} & 16 \\
\texttt{void*}               & 8  & 8  \\
\classe{std::string}         & 32 & 8  \\
\texttt{std::vector<int>}    & 24 & 8  \\
\classe{std::max\_align\_t}  & 16 & 16 \\
\bottomrule
\end{tabular}
\end{table}

Duas linhas merecem nota. O \texttt{long double} do x86-64 System V tem
80~\emph{bits} de dado útil e \texttt{sizeof}~16 --- quase metade é
\emph{padding} ---, e no MSVC ele é um \texttt{double} de 8~bytes, o que faz
dele o tipo menos portável do C++. E \classe{std::max\_align\_t} é o alinhamento
que \texttt{malloc} garante: 16 bytes. Pedir mais é o assunto seguinte.

\subsubsection{\texttt{alignas}: mandando no alinhamento}

\begin{porque}[três situações em que o alinhamento natural não basta]
\textbf{SIMD}: uma instrução AVX carrega 32~bytes de uma vez e exige o operando
alinhado a~32. \textbf{Cache}: dois campos escritos por \emph{threads}
diferentes na mesma linha de cache provocam \emph{false sharing} --- cada
escrita invalida a linha do outro núcleo, e o desempenho despenca sem que haja
\gls{datarace} nenhum. \textbf{Hardware}: um descritor de DMA que o dispositivo
exige em endereço específico.
\end{porque}

\begin{lstlisting}[style=bt, language=C++,
    caption={\texttt{alignas} em tipo e em membro, em
             \texttt{src/layout/demo\_alignas.hpp}.}]
struct alignas(32) Vetor8f { float v[8]; };     // pronto para AVX

struct Espacada {
    alignas(64) int produtor;                   // cada um na sua linha de cache
    alignas(64) int consumidor;
};
\end{lstlisting}

Medido pelo subprojeto: \texttt{Compacta} (os dois \texttt{int} juntos) ocupa
8~bytes; \texttt{Espacada} ocupa \textbf{128}, com o consumidor no
\emph{offset}~64. Dezesseis vezes mais memória para separar dois inteiros --- se
compensa é medição, não princípio, mas a alternativa pode custar uma ordem de
grandeza quando duas \emph{threads} os escrevem em laço.

\begin{armadilha}[\texttt{alignas} só aumenta; empacotar de verdade é fora do padrão]
\texttt{struct alignas(1) Tentativa \{ double d; \};} não compila ---
\texttt{alignment of 'Tentativa' is less than the minimum alignment (8) of its
members}. A recusa é proposital: ler um \texttt{double} desalinhado é \gls{ub}.
Para empacotar mesmo há extensões \textbf{não padronizadas}
(\texttt{\#pragma pack}, \texttt{\_\_attribute\_\_((packed))}), e elas criam um
perigo silencioso: um ponteiro para membro de \texttt{struct} empacotada é
\emph{desalinhado}, e passá-lo adiante é \gls{ub} --- o GCC avisa com
\texttt{-Waddress-of-packed-member}. Para protocolo de rede ou formato de
arquivo, serialize campo a campo.
\end{armadilha}

\paragraph{Alocação super-alinhada.} \texttt{malloc} garante 16 bytes; um tipo
que peça mais é \emph{over-aligned}. Antes do C++17 o \texttt{new}
\textbf{ignorava} o \texttt{alignas(32)} e devolvia memória alinhada a~16: o
código compilava, rodava e falhava na primeira instrução AVX. O C++17 (P0035)
acrescentou as sobrecargas com \classe{std::align\_val\_t}, e o compilador as
chama sozinho --- o subprojeto verifica que \textbf{8 de 8} alocações de
\texttt{Vetor8f} vieram alinhadas a 32. À mão há três formas:
\texttt{std::aligned\_alloc(32, n)}, \texttt{::operator new(n,
std::align\_val\_t\{32\})} e \classe{std::align}, que ajusta um bloco existente.

\begin{armadilha}[cada alocador tem o seu liberador, e \texttt{assume\_aligned} é promessa]
Memória de \texttt{aligned\_alloc} se libera com \texttt{free}; a de
\texttt{operator new(std::align\_val\_t)}, com o \texttt{operator delete}
correspondente --- misturar é \gls{ub}. Já
\texttt{std::assume\_aligned<32>(p)} (C++20) \emph{promete} ao compilador que o
ponteiro cru está alinhado, dispensando o prólogo escalar antes do laço
vetorizado. É promessa, não verificação: mentir funciona no x86 e falha no ARM.
O UBSan pega com \texttt{-fsanitize=alignment} (\cref{sec:sanitizers}).
\end{armadilha}

\subsubsection{As quatro categorias de tipo}
\label{sec:layout-categorias}

Quatro palavras parecidas, quatro autorizações diferentes. Não são curiosidade
acadêmica: cada uma \textbf{autoriza uma operação concreta}, e usar a operação
sem a autorização é \gls{ub}.

\begin{table}[htbp]
\centering
\caption{O que cada categoria de tipo autoriza.}
\label{tab:categorias-autoriza}
\footnotesize
\begin{tabular}{@{}ll@{}}
\toprule
\textbf{Categoria} & \textbf{Autoriza} \\
\midrule
\emph{trivially copyable} & \texttt{memcpy} do objeto inteiro --- é o que deixa
                            \texttt{vector<T>} usar \texttt{memmove} ao
                            realocar \\
\emph{standard-layout}    & \texttt{offsetof} legal; interoperação com C;
                            \texttt{\&s == \&s.primeiro\_membro} \\
\emph{agregado}           & inicialização por chaves e \emph{designated
                            initializers} \\
\emph{trivially destructible} & \texttt{vector<T>::clear()} só zera o tamanho,
                            sem percorrer \\
\bottomrule
\end{tabular}
\end{table}

\begin{lstlisting}[style=bt, language=bash,
    caption={Saída de \texttt{src/layout/}: os \emph{traits} instanciados sobre
             oito cobaias.}]
| tipo         | triv.cp | triv.ctor | std.layout | agregado |
|--------------|---------|-----------|------------|----------|
| int          |   sim   |    sim    |    sim     |   não    |
| Simples      |   sim   |    sim    |    sim     |   sim    |
| ComCtor      |   sim   |    não    |    sim     |   não    |
| ComVirtual   |   não   |    não    |    não     |   não    |
| Privado      |   sim   |    sim    |    não     |   não    |
| HerdaDados   |   sim   |    sim    |    não     |   sim    |
| ComString    |   não   |    não    |    sim     |   sim    |
| std::string  |   não   |    não    |    sim     |   não    |
\end{lstlisting}

Cada linha ensina uma regra. \texttt{ComCtor} (construtor escrito à mão) perdeu
\emph{trivially default constructible} e continua \emph{trivially copyable} ---
\texttt{memcpy} sobre ele é legal. \texttt{ComVirtual} perdeu tudo: o vptr é
preenchido pelo construtor, e um \texttt{memcpy} o corromperia.
\texttt{Privado}, com membros de dados em \texttt{public} \emph{e}
\texttt{private}, perdeu \emph{standard-layout}, porque aí o compilador pode
reordenar; \texttt{HerdaDados} perdeu pelo mesmo motivo, tendo dados em dois
níveis da hierarquia. E \texttt{ComString} não é trivial em nada e
\textbf{ainda é agregado}: \texttt{ComString\{1, "x"\}} compila.

\begin{armadilha}[\texttt{memcpy} sobre o tipo errado]
Sobre um \classe{std::string}, \texttt{memcpy} copia o \textbf{ponteiro
interno}: os dois objetos apontam para o mesmo \emph{buffer} e os dois
destrutores o liberam --- \emph{double free}. Com SSO e string curta pode até
``funcionar'', e é isso que torna o bug traiçoeiro: some nos testes. Sobre um
tipo com \texttt{virtual}, copia o \textbf{vptr}, e o destino passa a despachar
para os métodos da classe errada. A defesa cabe em duas linhas:
\begin{lstlisting}[style=bt, language=C++, frame=none,
                   xleftmargin=0pt, xrightmargin=0pt]
template <typename T>
void copia_bruta(T& d, const T& o) {
    static_assert(std::is_trivially_copyable_v<T>, "memcpy exige trivial");
    std::memcpy(&d, &o, sizeof(T));
}
\end{lstlisting}
Em C++20, \classe{std::bit\_cast} faz o mesmo com verificação em compilação e
sem \gls{ub} de \emph{aliasing} (\cref{sec:casts}).
\end{armadilha}

\paragraph{POD foi aposentado.} \texttt{std::is\_pod} está \textbf{depreciado
desde o C++20} (P0767). ``POD'' significava trivial \emph{e}
\emph{standard-layout}, e juntar as duas propriedades numa palavra só escondia
qual delas o programador precisava:

\begin{lstlisting}[style=bt, language=bash,
    caption={Saída de \texttt{src/layout/}: peça a autorização exata.}]
quero memcpy?          → is_trivially_copyable
quero offsetof?        → is_standard_layout
quero `T x;` de graça? → is_trivially_default_constructible
quero chaves?          → is_aggregate
\end{lstlisting}

As quatro perguntas são independentes, e \texttt{is\_pod} respondia ``sim'' só
quando \emph{todas} eram sim --- recusando tipos perfeitamente
\texttt{memcpy}-áveis, como o \texttt{ComCtor} da tabela.

\subsubsection{EBO e \texttt{[[no\_unique\_address]]}}
\label{sec:layout-ebo}

\begin{porque}[o problema: o 1 byte que vira 8]
Um objeto vazio ocupa 1~byte. Como \emph{membro}, esse byte custa muito mais: se
o membro seguinte é um \texttt{double}, ele só pode começar no \emph{offset}~8,
e os sete bytes intermediários viram \emph{padding}. Uma política vazia ---
comparador, alocador, \emph{deleter} --- custa oito bytes por objeto sem guardar
coisa nenhuma.
\end{porque}

\begin{lstlisting}[style=bt, language=C++,
    caption={As três formas, em
             \texttt{src/layout/demo\_no\_unique\_address.hpp}.}]
struct PorMembro {
    Vazio  politica;      // 1 byte...
    double dado;          // ...mas `dado` precisa de offset múltiplo de 8
};

struct PorHeranca : Vazio {
    double dado;
};

struct ComAtributo {
    [[no_unique_address]] Vazio politica;
    double                      dado;
};
\end{lstlisting}

\begin{lstlisting}[style=bt, language=bash,
    caption={Saída de \texttt{src/layout/}: metade da memória, num atributo.}]
struct { Vazio p; double d; }         → sizeof 16  ← 8 a mais!
struct : Vazio { double d; }          → sizeof  8  (EBO)
struct { [[nua]] Vazio p; double d; } → sizeof  8  (C++20)

Num std::vector de 10 milhões: 152.6 MB contra 76.3 MB
\end{lstlisting}

A saída antiga chama-se \textbf{EBO} --- \emph{Empty Base Optimization}: uma
classe base vazia pode ocupar zero bytes no derivado, porque a linguagem não
exige endereço próprio para subobjetos de base. É por isso que a libstdc++ faz o
\texttt{std::vector} \textbf{herdar} do alocador:

\begin{lstlisting}[style=bt, language=bash,
    caption={Saída de \texttt{src/layout/}: a EBO na biblioteca padrão.}]
sizeof(std::allocator<int>)         =  1 (é VAZIO)
sizeof(std::vector<int>)            = 24 (3 ponteiros, nada mais)
→ se o alocador fosse membro, o vector teria 32.

unique_ptr<int>                     =  8 bytes
unique_ptr<int, lambda SEM captura> =  8 bytes (EBO)
unique_ptr<int, void(*)(int*)>      = 16 bytes (ponteiro!)
\end{lstlisting}

\begin{armadilha}[passar \texttt{\&funcao} como \emph{deleter} dobra o \texttt{unique\_ptr}]
Um \emph{lambda} sem captura é um tipo vazio e some pela EBO; um ponteiro de
função é um objeto de 8~bytes que precisa ser guardado. O
\classe{std::unique\_ptr} sai de 8 para 16 bytes, deixa de caber num registrador
e dobra o custo de guardá-lo em container (\cref{sec:smart-pointers}).
\end{armadilha}

Herdar de algo que não é interface incomoda, e era a única saída antes do C++20
--- o \texttt{boost::compressed\_pair} existia para embrulhar essa feiura. O
\texttt{[[no\_unique\_address]]} diz o mesmo sem herança: ``este membro pode
compartilhar endereço com outro''.

\begin{armadilha}[dois membros vazios do mesmo tipo não colapsam]
Dois subobjetos do \textbf{mesmo tipo} não podem ter o mesmo endereço, e a regra
vale mesmo com o atributo. Mas isso não significa que o custo apareça na
\texttt{struct} que os contém: basta haver outro membro para absorvê-los.
\begin{lstlisting}[style=bt, language=bash, frame=none,
                   xleftmargin=0pt, xrightmargin=0pt]
struct { [[nua]] Vazio a; [[nua]] Vazio b;      } → sizeof 2
struct { [[nua]] Vazio a; [[nua]] Comparador b; } → sizeof 1
struct { [[nua]] Vazio a, b; int dado;          } → sizeof 4
  offset de a = 0   offset de b = 1   offset de dado = 0
\end{lstlisting}
Os dois vazios têm endereços diferentes entre si, como manda a regra, e ao mesmo
tempo moram \emph{dentro} do \texttt{int} --- que é exatamente o que o atributo
autoriza.
\end{armadilha}

\begin{armadilha}[o atributo é autorização, não garantia]
O compilador pode ignorá-lo, e o MSVC de fato o ignora por compatibilidade de
\gls{abi}, oferecendo \texttt{[[msvc::no\_unique\_address]]} no lugar. Se o
\emph{layout} importa, trave o contrato com
\texttt{static\_assert(sizeof(T)~==~N)} em vez de confiar.
\end{armadilha}

\begin{regra}[toda política passada por \emph{template} merece o atributo]
Comparador, alocador, \emph{deleter}, \emph{hasher}, política de \emph{mutex}:
quase sempre vazios, e todos custam 8 bytes por objeto se forem membros comuns.
O subprojeto mede o caso num \texttt{ConjuntoOrdenado<T, Comp>}: o
\texttt{std::less<int>} custou \textbf{0 bytes}; um comparador com um
\texttt{int} de estado custou~8 --- e aí é justo pagar.
\end{regra}

\subsubsection{O vptr, e o que \texttt{virtual} custa em bytes}
\label{sec:layout-vptr}

A \cref{sec:oop} tratou do \emph{comportamento} do despacho dinâmico; aqui o
assunto é o preço em bytes e em endereços. Quando uma classe tem ao menos uma
função \texttt{virtual}, o compilador cria para ela uma \textbf{vtable} --- um
vetor de ponteiros de função, guardado \emph{no binário}, um por classe --- e
acrescenta a cada \emph{objeto} um \textbf{vptr}, ponteiro para a vtable da sua
classe real. É o vptr que permite a \texttt{forma->area()} chamar a
implementação certa sem saber o tipo.

\begin{lstlisting}[style=bt, language=bash,
    caption={Saída de \texttt{src/layout/}: o custo de um único
             \texttt{virtual}.}]
struct { int a; int b; }                  → sizeof  8
struct { int a; int b; virtual ~T(); }    → sizeof 16

os 16 bytes de um ComVirtual{a=0x11111111, b=0x22222222}:
  80 E7 47 87 CA 55 00 00 | 11 11 11 11 22 22 22 22
  └─── vptr (8 bytes) ───   └─── a ───   └─── b ───
offsetof(ComVirtual, a) = 8 — o vptr vem ANTES.
\end{lstlisting}

Um \emph{único} método virtual custa 8 bytes por objeto; dez métodos custam os
\emph{mesmos} 8, porque o vptr é um só e a vtable, que cresce, fica no binário.
A posição no início é convenção do \textbf{Itanium C++ ABI} (GCC e Clang em
Linux e macOS), não exigência do padrão --- mas é o que permite despachar sem
conhecer o tipo, já que o vptr está sempre no \emph{offset}~0.

\begin{regra}[destrutor virtual só onde há destruição por ponteiro-para-base]
Num \texttt{vector<Particula>} de dez milhões, um
\texttt{virtual~\textasciitilde Particula()} que ninguém usa custa
\textbf{76~MB}. Se a classe não é destruída através de ponteiro para a base, não
declare o destrutor virtual: torne-o \texttt{protected} e não virtual, o que
impede o uso errado sem cobrar o vptr.
\end{regra}

\paragraph{Herança simples, e uma surpresa.} Com herança simples o subobjeto da
base fica no início, o derivado acrescenta depois e o \emph{upcast} é gratuito.
Mas a medição contraria a expectativa:

\begin{lstlisting}[style=bt, language=bash,
    caption={Saída de \texttt{src/layout/}: o derivado tem um \texttt{int} a
             mais e o mesmo \texttt{sizeof}.}]
sizeof(Base1)  = 16   sizeof(Deriv1) = 16
Base1:  [vptr 8][x 4][padding 4]
Deriv1: [vptr 8][x 4][   y   4]     ← y foi para DENTRO do padding de cauda
offsetof: x=8, y=12
\end{lstlisting}

O compilador só pode reaproveitar o \emph{padding} de cauda da base porque
\texttt{Base1} \textbf{não é} \emph{standard-layout} (ela tem \texttt{virtual});
se fosse, o padrão proibiria. Aqui a \cref{sec:layout-categorias} encontra a
prática: é por isso que \texttt{memcpy} de um objeto para o
\emph{subobjeto-base} de outro pode corromper o derivado --- os bytes de
``\emph{padding}'' da base, no destino, são o \texttt{y} de alguém. E é esse
mesmo \emph{layout} que torna o \gls{slicing} possível: copiar \texttt{*pb} para
um \texttt{Base1} leva só a parte-base, e o resultado é um \texttt{Base1}
legítimo, sem aviso nenhum.

\paragraph{Herança múltipla: o \texttt{this} que se desloca.} Com duas bases
polimórficas o objeto tem \textbf{dois} vptr --- a parte-B precisa da sua
própria vtable, e um \texttt{B*} tem de apontar para o começo dela.

\begin{lstlisting}[style=bt, language=C++,
    caption={As três classes de
             \texttt{src/layout/demo\_polimorfico.hpp}.}]
struct A { int a = 1; virtual ~A() = default; virtual void fa() {} };
struct B { int b = 2; virtual ~B() = default; virtual void fb() {} };
struct AB : A, B { int c = 3; };
\end{lstlisting}

\begin{lstlisting}[style=bt, language=bash,
    caption={Saída de \texttt{src/layout/}: o mesmo objeto, dois endereços.}]
&obj       offset  0
(A*)&obj   offset  0   ← igual ao início
(B*)&obj   offset 16   ← DESLOCADO!
Layout: [vptr-A][a][pad][vptr-B][b][pad][c][pad]

(void*)pa == (void*)pb ? false
pa == &obj && pb == &obj ? true
dynamic_cast<void*>(pa) == dynamic_cast<void*>(pb) ? true
\end{lstlisting}

\begin{armadilha}[três consequências da herança múltipla]
\textbf{(1)}~\texttt{static\_cast<B*>(\&obj)} \textbf{não é um no-op}: o
compilador soma~16. Um \classe{reinterpret\_cast} no lugar produziria um
ponteiro errado que compila sem aviso --- é a diferença entre os dois
\emph{casts} da \cref{sec:casts}, tornada visível.
\textbf{(2)}~Comparar ponteiros engana: \texttt{(void*)pa == (void*)pb} é
\emph{false} embora os dois apontem para o mesmo objeto; para comparar
identidade, use \texttt{dynamic\_cast<void*>}, que devolve o início do objeto
mais derivado. \textbf{(3)}~\texttt{delete pb} precisa \emph{ajustar} o ponteiro
antes de liberar, e quem faz isso é um \emph{thunk} na vtable de~B --- sem
destrutor virtual em~B, seria \gls{ub}.
\end{armadilha}

\paragraph{Herança virtual: o diamante que custa mais tendo menos.} Quando duas
classes herdam da mesma base e uma terceira herda das duas, ou a base aparece
duas vezes, ou é compartilhada; \texttt{virtual} escolhe a segunda opção. O
número contraria a expectativa: \texttt{sizeof(NaoDiamante)}~=~\textbf{40}, com
a base~\texttt{V} duplicada, contra \texttt{sizeof(Diamante)}~=~\textbf{48}, com
ela compartilhada. O virtual custou \emph{mais} bytes tendo \emph{menos} dados.

O mecanismo explica o paradoxo. \texttt{E1} e \texttt{E2} não sabem, em
compilação, \emph{onde} \texttt{V} vai estar --- quem decide é a classe mais
derivada, que pode nem existir ainda. Cada uma ganha então um ponteiro extra, o
\emph{virtual base offset}, consultado em tempo de execução, e o \texttt{V}
compartilhado vai para o fim do objeto. Acessar \texttt{d.v} passa a custar uma
indireção a mais que acessar \texttt{d.d}, e construir \texttt{V} vira
responsabilidade da classe mais derivada, pulando os construtores de \texttt{E1}
e \texttt{E2}.

\begin{regra}[herança virtual só para interfaces sem dados]
Com bases puramente abstratas o custo é pequeno e o diamante é inevitável. Com
\emph{dados} na base virtual, reconsidere o desenho: composição costuma ser mais
simples, mais rápida e mais fácil de explicar.
\end{regra}

\subsubsection{Não adivinhe: inspecione}

\begin{lstlisting}[style=bt, language=bash,
    caption={As quatro ferramentas de inspeção de \emph{layout}.}]
clang++ -Xclang -fdump-record-layouts    # layout completo, clang
g++     -fdump-lang-class                # gera arquivos .class
pahole -C NomeDaStruct ./binario         # mostra os BURACOS
static_assert(sizeof(T) == 16);          # trava o contrato
\end{lstlisting}

O \texttt{static\_assert} é o mais valioso e o único que sobrevive ao tempo:
transforma uma suposição sobre \emph{layout} numa falha de compilação no dia em
que alguém acrescentar um membro. O \texttt{pahole}, do pacote
\texttt{dwarves}, lê a informação de depuração do binário e imprime literalmente
os buracos de \emph{padding} de cada \texttt{struct}.

Vale fechar dizendo quando nada disso importa: com dezenas ou centenas de
objetos, com um objeto atrás de um ponteiro seguido uma vez por quadro, ou
quando a clareza vale mais que 8~bytes. O \emph{layout} importa quando há
milhões de objetos, quando os dados atravessam uma fronteira de \gls{abi} ou vão
para disco, quando há interoperação com C ou GPU, e quando o laço quente
percorre um \emph{array} e o gargalo é o cache.

% =============================================================================
\subsection{Apagamento de tipo}
\label{sec:erasure}
% =============================================================================

\subsubsection{O problema, e o que a herança cobra por resolvê-lo}

\begin{porque}[o buraco no sistema de tipos estático]
Três \texttt{struct}s que não se conhecem e não herdam de nada ---
\texttt{Circulo}, \texttt{Quadrado}, \texttt{Triangulo}, cada uma com
\texttt{area()}. Como guardar as três num único \texttt{std::vector}? Não há
resposta boa: \texttt{vector<Circulo>} só guarda círculos; um \emph{template} é
resolvido em compilação e não dá heterogeneidade; e \texttt{void*} perde toda a
informação, inclusive \emph{como destruir}.
\end{porque}

\Gls{typeerasure} --- apagamento de tipo --- é guardar um objeto de tipo
\texttt{T} num lugar que \textbf{não menciona \texttt{T}}, preservando as
operações que interessam. O tipo some da assinatura; o comportamento fica. Você
já usa isso sem chamar pelo nome:

\begin{lstlisting}[style=bt, language=C++,
    caption={Quatro apagadores de tipo da biblioteca padrão.}]
std::function<int(int)> f;   // aceita lambda, functor, ponteiro de função...
std::any                a;   // aceita qualquer coisa copiável
std::shared_ptr<void>   p;   // guarda o deleter de QUALQUER tipo
std::unique_ptr<Base>   b;   // a herança é a forma mais antiga
\end{lstlisting}

A última linha é a que quase ninguém percebe.
\texttt{std::unique\_ptr<Forma>} não menciona \texttt{Circulo} nem
\texttt{Quadrado}: o tipo foi apagado, e a \gls{vtable} é o mecanismo. Herança
polimórfica \emph{é} apagamento de tipo --- só que \textbf{intrusivo}: para
entrar no vetor, o tipo precisa \emph{herdar}. Isso exclui \texttt{int},
\classe{std::string} e qualquer tipo de terceiros que você não pode alterar;
acrescenta um vptr ao seu tipo, que deixa de ser trivial e de caber num
registrador (\cref{sec:layout-vptr}); obriga a alocar no \emph{heap} e usar
ponteiro, sob pena de \gls{slicing}; e acopla o seu tipo a uma hierarquia que
ele não escolheu.

\begin{table}[htbp]
\centering
\caption{As três formas de guardar coisas heterogêneas.}
\label{tab:formas-heterogeneas}
\footnotesize
\begin{tabular}{@{}lll@{}}
\toprule
\textbf{Forma} & \textbf{Conjunto de tipos} & \textbf{Despacho} \\
\midrule
herança + \texttt{virtual}             & \textbf{aberto}, intrusivo & vtable \\
\classe{std::variant} + \texttt{visit} & \textbf{fechado}           & índice / \emph{jump table} \\
apagamento de tipo manual              & \textbf{aberto}, não intrusivo & vtable manual \\
\bottomrule
\end{tabular}
\end{table}

\emph{Aberto} significa que qualquer um pode acrescentar um tipo novo depois, sem
tocar no seu código; \emph{fechado}, que a lista está fixa na declaração. A
terceira linha é o assunto desta seção.

\subsubsection{\texttt{std::function} do zero: Conceito, Modelo, Wrapper}

A receita tem sempre as mesmas três peças (\cref{fig:erasure-pecas}): um
\textbf{Conceito}, interface abstrata que descreve as operações mas não o tipo;
um \textbf{Modelo}, \emph{template} sobre \texttt{T}, que implementa a interface
delegando para o \texttt{T}; e um \textbf{Wrapper} com semântica de valor, que
guarda um ponteiro para o Conceito.

\begin{figure}[htbp]
\centering
\begin{tikzpicture}[node distance=7mm and 10mm]
  \node[caixaforte] (w) {\textbf{Wrapper}\\\scriptsize \texttt{Funcao<R(Args...)>}\\\scriptsize valor, copiável};
  \node[caixa, below=8mm of w] (c) {\textbf{Conceito}\\\scriptsize interface abstrata\\\scriptsize \texttt{invocar} / \texttt{clonar}};
  \node[caixa, below left=9mm and 6mm of c]  (m1) {\textbf{Modelo}\texttt{<Triplica>}};
  \node[caixa, below right=9mm and 6mm of c] (m2) {\textbf{Modelo}\texttt{<lambda>}};
  \node[folha, below=7mm of m1] (t1) {\texttt{Triplica}\\\scriptsize não herda de nada};
  \node[folha, below=7mm of m2] (t2) {\emph{lambda}\\\scriptsize não herda de nada};
  \draw[seta] (w) -- node[rotulo,right] {\texttt{unique\_ptr}} (c);
  \draw[aresta] (c) -- (m1); \draw[aresta] (c) -- (m2);
  \draw[seta] (m1) -- node[rotulo,left] {guarda} (t1);
  \draw[seta] (m2) -- node[rotulo,right] {guarda} (t2);
\end{tikzpicture}
\caption{As três peças do apagamento de tipo. Só o \textbf{Modelo} conhece o
         \texttt{T}: ele herda do Conceito e some atrás do ponteiro. Quem
         escreveu o \texttt{T} nunca ouviu falar de nenhuma das três.}
\label{fig:erasure-pecas}
\end{figure}

\begin{lstlisting}[style=bt, language=C++,
    caption={A receita completa (resumo de
             \texttt{src/type\_erasure/demo\_implementacao.hpp}).}]
template <typename R, typename... Args>
class Funcao<R(Args...)> {
    // PEÇA 1: o Conceito -- o que sabemos fazer, sem saber com o quê
    struct Conceito {
        virtual ~Conceito() = default;
        virtual R                          invocar(Args...) const = 0;
        virtual std::unique_ptr<Conceito>  clonar()          const = 0;
    };

    // PEÇA 2: o Modelo -- conhece T, implementa o Conceito
    template <typename T>
    struct Modelo final : Conceito {
        T alvo;
        explicit Modelo(T t) : alvo{std::move(t)} {}

        R invocar(Args... a) const override {
            return const_cast<T&>(alvo)(std::forward<Args>(a)...);
        }
        std::unique_ptr<Conceito> clonar() const override {
            return std::make_unique<Modelo>(alvo);   // exige T copiável
        }
    };

    // PEÇA 3: o Wrapper -- valor semântico
    std::unique_ptr<Conceito> pimpl_;

public:
    template <typename T>
        requires (!std::is_same_v<std::decay_t<T>, Funcao>)
    Funcao(T&& t)
        : pimpl_{std::make_unique<Modelo<std::decay_t<T>>>(std::forward<T>(t))} {}
};
\end{lstlisting}

Com isso, um \emph{functor}, um \emph{lambda} com captura e um ponteiro de
função --- três tipos sem base comum, que não sabem que \texttt{Funcao} existe
--- entram no mesmo objeto. O acoplamento é zero: o único contrato é ser
invocável com \texttt{(Args...)} devolvendo algo conversível para \texttt{R}.

\begin{armadilha}[sem o \texttt{requires}, o construtor \emph{template} sequestra a cópia]
O construtor \emph{template} aceita qualquer coisa --- inclusive uma
\texttt{Funcao}. Sem a restrição, \texttt{Funcao f2\{f1\}} escolheria o
\emph{template} em vez do construtor de cópia e construiria uma
\texttt{Funcao<Funcao<...>>}: um objeto que embrulha o outro em vez de copiá-lo.
Compila, roda e faz outra coisa. Toda classe com construtor \emph{template}
irrestrito tem esse bug latente.
\end{armadilha}

\paragraph{Por que \texttt{clonar()}.} \texttt{unique\_ptr<Conceito>} não é
copiável --- logo, sem o \texttt{clonar()} virtual, a \texttt{Funcao} deixaria
de ser um \emph{valor}. Há exatamente três saídas, e cada biblioteca escolheu
uma: \texttt{clonar()} virtual dá cópia profunda (é o padrão \emph{Prototype}, e
o que a \classe{std::function} faz); \texttt{shared\_ptr<const Conceito>} dá
cópia rasa e barata, com o alvo compartilhado; proibir a cópia é o
\classe{std::move\_only\_function} do C++23. A \texttt{std::function} escolheu a
primeira e paga o preço: ela \textbf{exige} \emph{callable} copiável, e um
\emph{lambda} que captura um \texttt{unique\_ptr} por \gls{move} não cabe nela
--- \texttt{error: static assertion failed: std::function target must be
copy-constructible}.

\subsubsection{SBO: medido, não presumido}

\begin{porque}[o problema que o \emph{buffer} pequeno resolve]
A \texttt{Funcao} acima aloca no \emph{heap} \textbf{sempre}, mesmo para um
\emph{lambda} sem captura --- que não guarda nada. Como o \emph{callback}
pequeno é o caso dominante, as implementações reais reservam alguns bytes
\emph{dentro} do próprio objeto e constroem o alvo ali quando ele cabe: é a
\textbf{SBO}, \emph{small buffer optimization}. Cabe até quanto? O padrão não
diz, então só resta medir.
\end{porque}

O subprojeto substitui o \texttt{operator new} global por uma versão que conta
alocações e bytes, e constrói \emph{lambdas} com capturas de tamanho conhecido:

\begin{lstlisting}[style=bt, language=bash,
    caption={Saída de \texttt{src/type\_erasure/}: onde exatamente o
             \texttt{new} aparece.}]
sizeof(std::function<void()>)  = 32 bytes
sizeof(v1::Funcao<void()>)     =  8 bytes  (só o unique_ptr)

std::function, sem captura → 0 alocação(ões), INLINE no buffer
std::function, captura 8B  → 0 alocação(ões), INLINE no buffer
std::function, captura 16B → 0 alocação(ões), INLINE no buffer
std::function, captura 24B → 1 alocação(ões), 24 bytes
std::function, captura 64B → 1 alocação(ões), 64 bytes
v1::Funcao, sem captura    → 1 alocação(ões), SEMPRE
\end{lstlisting}

O corte está entre 16 e 24 bytes, e o \texttt{sizeof} de 32 explica por quê:
ponteiro de vtable, ponteiro de invocação e 16 bytes de \emph{buffer}. A regra é
da \textbf{implementação}, não do padrão --- e a libstdc++ acrescenta uma
condição própria: o \emph{callable} precisa ser
\emph{nothrow-move-constructible} para usar o \emph{buffer}; se o \gls{move}
puder lançar, ela aloca, porque mover a \texttt{function} teria de mover o
\emph{buffer}. É o mesmo raciocínio que faz o \texttt{std::vector} copiar em vez
de mover ao realocar (\cref{sec:value-semantics}).

\begin{regra}[o que cabe de graça numa \texttt{std::function}]
Um \emph{lambda} que captura \texttt{[this]}, um \texttt{int} ou dois ponteiros
é gratuito. Um que captura três \classe{std::string} por valor aloca --- na
construção e em cada cópia. Em código quente, prefira passar o \emph{callable}
por \emph{template} e reserve a \texttt{std::function} para as fronteiras em que
o tipo precisa mesmo ser apagado.
\end{regra}

\begin{armadilha}[substituir o \texttt{operator new} global tem três regras não óbvias]
\textbf{(1)}~A substituição precisa estar no \textbf{escopo global}:
declará-la dentro de um \emph{namespace} é erro de compilação, porque ali ela
não seria substituição, seria outra função com o mesmo nome. \textbf{(2)}~Precisa
ser \textbf{única no programa} --- é uma função com \emph{linkage} externa
reservada ---, e por isso não pode ser \texttt{inline}: o cabeçalho que a contém
só pode ser incluído por \emph{uma} unidade de tradução (\cref{sec:linkage}).
\textbf{(3)}~Os \texttt{operator delete} pareados também precisam ser
substituídos: misturar o \texttt{new} substituído com o \texttt{delete} padrão é
\gls{ub}. De quebra, o GCC emite um \texttt{-Wmismatched-new-delete}
falso-positivo ao \emph{inlinar} os corpos, suprimido no subprojeto com
\emph{pragma} local.
\end{armadilha}

\subsubsection{Vtable manual, sem herança nenhuma}

A versão acima usa \texttt{virtual}, o que é legível e caro: o objeto guardado
ganha um vptr, e vptr no \emph{buffer} complica o \gls{move}. Dá para apagar o
tipo \textbf{sem herança}: um \texttt{struct} de ponteiros de função,
\texttt{static constexpr}, uma instância por tipo \texttt{T}.

\begin{lstlisting}[style=bt, language=C++,
    caption={A vtable escrita à mão (resumo de
             \texttt{src/type\_erasure/demo\_implementacao.hpp}).}]
template <typename R, typename... Args>
class Funcao {
    // A "vtable": um struct de ponteiros de função, UMA instância estática
    // por tipo T. Sem `virtual`, sem vptr no objeto.
    struct VTable {
        R    (*invocar)(void*, Args...);
        void (*destruir)(void*);
        void*(*clonar)(void*);
    };

    template <typename T>
    static const VTable* tabela_de() {
        static constexpr VTable t{
            +[](void* p, Args... a) -> R {
                return (*static_cast<T*>(p))(std::forward<Args>(a)...);
            },
            +[](void* p) { delete static_cast<T*>(p); },
            +[](void* p) -> void* { return new T{*static_cast<T*>(p)}; },
        };
        return &t;
    }

    void*         alvo_ = nullptr;
    const VTable* vt_   = nullptr;
};
\end{lstlisting}

É literalmente o que o compilador faz com \texttt{virtual}, escrito à mão --- e
essa diferença aparentemente estética muda o que é possível
(\cref{tab:vtable-manual}). Sem vptr, o objeto guardado pode ir direto para o
\emph{buffer} interno: é por isso que a libstdc++ escolhe esta forma. O
\texttt{virtual} da primeira versão é mais claro e serve para aprender.

\begin{table}[htbp]
\centering
\caption{\texttt{virtual} contra vtable manual.}
\label{tab:vtable-manual}
\footnotesize
\begin{tabular}{@{}lll@{}}
\toprule
\textbf{Aspecto} & \texttt{virtual} & \textbf{vtable manual} \\
\midrule
ponteiro no \textbf{objeto}   & sim (vptr) & \textbf{não} \\
permite SBO                   & difícil    & \textbf{sim} \\
escolher as operações         & fixo       & por chamada \\
legibilidade                  & melhor     & pior \\
\bottomrule
\end{tabular}
\end{table}

\subsubsection{O apagamento que já vem no padrão}

Cada apagador da biblioteca padrão preserva um conjunto diferente de operações,
e é essa escolha que os distingue: \classe{std::function} preserva
\textbf{chamar}; \classe{std::any}, destruir e copiar (nem imprimir dá);
\texttt{std::shared\_ptr<T>}, \textbf{destruir certo};
\classe{std::move\_only\_function}, chamar sem exigir cópia;
\classe{std::type\_index}, identidade de tipo em execução.

\texttt{std::any} exige o tipo \textbf{exato} na recuperação:
\texttt{any\_cast<long>} sobre um \texttt{int} lança \texttt{bad\_any\_cast},
porque a comparação é \texttt{typeid} exato e não há conversão nenhuma --- a
forma de ponteiro devolve \texttt{nullptr} em vez de lançar. Quando usar: quase
nunca. Antes, veja se um \texttt{variant} serve (\cref{sec:vocabulary}) ou se um
apagamento sob medida preserva exatamente as operações necessárias. \texttt{any}
é para guardar algo e devolvê-lo intacto \textbf{sem nunca olhar dentro}.

\begin{armadilha}[o \emph{deleter} apagado do \texttt{shared\_ptr}]
Em \texttt{std::unique\_ptr<T, D>} o \emph{deleter} está \textbf{no tipo}; em
\texttt{std::shared\_ptr<T>}, está no \textbf{bloco de controle}, com o tipo
apagado --- e isso tem três consequências. \textbf{(1)}~Um
\texttt{shared\_ptr<Base>} destrói corretamente mesmo \emph{sem} destrutor
virtual, porque o \emph{deleter} foi capturado quando o tipo ainda era o
derivado (o subprojeto verifica que \texttt{\textasciitilde Contada} roda). Isso
\textbf{não} é licença para omitir o destrutor virtual: só funciona se o
\texttt{shared\_ptr} foi criado a partir do tipo derivado.
\textbf{(2)}~Recursos de C saem de graça:
\texttt{std::shared\_ptr<FILE> f\{fopen(...), \&fclose\}}.
\textbf{(3)}~O \emph{aliasing constructor} guarda um ponteiro e mantém vivo
\emph{outro} objeto. O preço é o bloco de controle --- uma alocação, evitável
com \texttt{make\_shared} --- e uma indireção a mais na destruição.
\end{armadilha}

O \texttt{std::move\_only\_function} do C++23 aceita o \emph{lambda} que captura
por \gls{move} e, de quebra, acertou os qualificadores da assinatura:
\texttt{move\_only\_function<int() const>} tem \texttt{operator()} constante de
verdade, \texttt{<int() \&\&>} só é invocável em \emph{rvalue}, e
\texttt{<int() noexcept>} propaga o \texttt{noexcept}. A
\texttt{std::function} faz um \classe{const\_cast} interno --- exatamente como a
versão didática acima --- e portanto \emph{mente sobre const-correção}.

Já o \texttt{std::type\_index} é o embrulho que torna \texttt{std::type\_info}
--- não copiável e sem \texttt{operator<} --- utilizável como chave de
\texttt{map}. Com ele se constrói o despacho dinâmico \textbf{aberto}: um
\texttt{map<type\_index, handler>} em que qualquer módulo registra um tipo novo
sem tocar em nada, base do \emph{plugin registry}
(\cref{sec:plugin-registry}) e do ECS (\cref{sec:ecs}). O custo é uma busca em
\texttt{std::map} por chamada, contra uma indireção só na vtable.

\begin{armadilha}[\texttt{typeid(T).name()} é manglado e não é portável]
O nome devolvido é o símbolo decorado pelo compilador e muda entre compiladores
e entre versões do mesmo compilador. Para exibir, use
\texttt{abi::\_\_cxa\_demangle} do GCC e do Clang. \textbf{Nunca serialize esse
nome}: um arquivo gravado hoje pode não ser lido pelo binário de amanhã.
\end{armadilha}

\subsubsection{A medição que decide}
\label{sec:erasure-medicao}

Tabelas comparando \texttt{virtual}, \gls{crtp}, \texttt{variant} e apagamento
de tipo existem em todo lugar; o que falta é o número. O subprojeto soma a área
de 200\,000 formas, alternando círculos e quadrados para que a predição de
desvio não favoreça ninguém, com média de 20 rodadas e \texttt{-O2}:

\begin{lstlisting}[style=bt, language=bash,
    caption={Saída de \texttt{src/type\_erasure/}: 200\,000 elementos, 20
             rodadas, \texttt{-O2}.}]
virtual (unique_ptr<Base>)           0.514 ms/rodada    3.9×
CRTP (vetores homogêneos)            0.130 ms/rodada    1.0×
variant + visit                      0.131 ms/rodada    1.0×
type erasure (unique_ptr interno)    0.517 ms/rodada    4.0×

Memória por elemento:
  virtual      :  8 B (ponteiro) + heap: vptr + dado
  CRTP         :  8 B, contíguo, SEM vptr
  variant      : 16 B, contíguo (maior alternativa + tag)
  type erasure :  8 B (ponteiro) + heap
\end{lstlisting}

\begin{porque}[o que domina o resultado não é o despacho]
A leitura ingênua atribuiria o 4× ao custo da chamada indireta, e os próprios
números a desmentem: \texttt{variant + visit} \emph{também} é despacho dinâmico
--- um salto por tabela --- e empatou com o CRTP, que não tem despacho nenhum. O
que separa os dois grupos é a \textbf{localidade}: \texttt{variant} e CRTP
percorrem memória contígua, enquanto \texttt{virtual} e apagamento de tipo
percorrem um vetor de \emph{ponteiros} para objetos espalhados no \emph{heap}.
Uma falta de cache custa cerca de 200 ciclos; uma chamada indireta bem predita,
cerca de~2. O 4× é o cache falando, não a vtable --- e a consequência prática é
que eliminar o \texttt{virtual} sem mudar o \emph{layout} não resolveria nada.
\end{porque}

\subsubsection{O problema da expressão}

Você tem \textbf{tipos} (círculo, quadrado, \ldots) e \textbf{operações} (área,
perímetro, desenhar, \ldots), e só dá para deixar \emph{um} dos dois eixos
barato de estender.

\begin{table}[htbp]
\centering
\caption{O problema da expressão: cada abordagem escolhe um eixo.}
\label{tab:expressao}
\footnotesize
\begin{tabular}{@{}lll@{}}
\toprule
\textbf{Abordagem} & \textbf{Tipo novo} & \textbf{Operação nova} \\
\midrule
\texttt{virtual}         & fácil: 1 classe & \textbf{caro}: mexe na base e em
                                             todo derivado \\
\texttt{variant + visit} & \textbf{caro}: recompila todos os \emph{visitors}
                                           & fácil: outro \emph{visitor} \\
apagamento de tipo       & fácil, e sem tocar no tipo & \textbf{caro}: mexe no
                                             \texttt{Conceito} \\
\gls{crtp}               & fácil & fácil --- \emph{mas sem heterogeneidade em
                                   execução} \\
\bottomrule
\end{tabular}
\end{table}

A última linha é uma armadilha de leitura: o CRTP parece ganhar dos dois lados
porque \textbf{não resolve o mesmo problema} --- ele não dá container
heterogêneo. Se você precisa de um \texttt{vector<Forma>} com tipos decididos em
execução, o CRTP não é candidato; foi por isso que a medição da
\cref{sec:erasure-medicao} precisou dar a ele \emph{dois vetores homogêneos}.

\begin{regra}[como decidir, na ordem]
\textbf{1.}~Tipos conhecidos, lista curta e estável? $\rightarrow$
\texttt{variant}: contíguo, sem alocação, exaustividade verificada em
compilação. \textbf{2.}~Tipos abertos e você controla as classes? $\rightarrow$
herança + \texttt{virtual}, o mais legível. \textbf{3.}~Tipos abertos e você
\textbf{não} controla as classes (\texttt{int}, \texttt{std::string},
terceiros)? $\rightarrow$ apagamento de tipo. \textbf{4.}~Sem heterogeneidade em
execução, só reuso de código? $\rightarrow$ \emph{template}/CRTP, e nem pense em
vtable.
\end{regra}

\subsubsection{Polimorfismo com semântica de valor}

A crítica de Sean Parent (\emph{Inheritance Is The Base Class of Evil},
GoingNative 2013) cabe em uma frase: herdar para ser polimórfico contamina o seu
tipo com ponteiro, \emph{heap} e \emph{aliasing} --- problemas que ele não
tinha. A alternativa é uma inversão simples de enunciar e difícil de perceber
sozinho: \textbf{o polimorfismo mora no container, não no tipo guardado}.

\begin{lstlisting}[style=bt, language=C++,
    caption={O tipo-valor polimórfico (resumo de
             \texttt{src/type\_erasure/demo\_pratica.hpp}).}]
// O contrato, explícito. Sem ele, um tipo inadequado gera erro dentro de
// Modelo<T> -- dez níveis de template abaixo da linha do usuário.
template <typename T>
concept SabeDesenhar = requires(const T& t, int nivel) {
    { t.desenhar(nivel) } -> std::same_as<void>;
    { t.nome() }          -> std::convertible_to<std::string>;
};

class Desenhavel {
    struct Conceito {
        virtual ~Conceito() = default;
        virtual void                      desenhar(int) const = 0;
        virtual std::unique_ptr<Conceito> clonar()      const = 0;
    };
    template <typename T>
    struct Modelo final : Conceito {
        T v;
        explicit Modelo(T x) : v{std::move(x)} {}
        void desenhar(int n) const override { v.desenhar(n); }
        std::unique_ptr<Conceito> clonar() const override {
            return std::make_unique<Modelo>(v);
        }
    };
    std::unique_ptr<Conceito> p_;

public:
    template <SabeDesenhar T>
    Desenhavel(T x) : p_{std::make_unique<Modelo<T>>(std::move(x))} {}

    // Rule of 5 -- é isto que dá semântica de valor
    Desenhavel(const Desenhavel& o) : p_{o.p_->clonar()} {}
    Desenhavel& operator=(const Desenhavel& o) {
        if (this != &o) p_ = o.p_->clonar();
        return *this;
    }
    Desenhavel(Desenhavel&&) noexcept            = default;
    Desenhavel& operator=(Desenhavel&&) noexcept = default;
    ~Desenhavel()                                = default;

    void desenhar(int n = 0) const { p_->desenhar(n); }
};
\end{lstlisting}

O que isso compra é uma cena copiável:

\begin{lstlisting}[style=bt, language=C++,
    caption={Uma cena por valor, com cópia profunda.}]
std::vector<Desenhavel> cena;          // por VALOR, não unique_ptr
cena.emplace_back(Ponto{1, 2});
cena.emplace_back(Texto{"olá"});
cena.emplace_back(Grupo{"sub", {Ponto{9,9}, Texto{"aninhado"}}});

std::vector<Desenhavel> copia = cena;  // PROFUNDA -- nem compilaria
                                       // com unique_ptr<Base>
\end{lstlisting}

\texttt{Ponto}, \texttt{Texto} e \texttt{Grupo} são \texttt{struct}s comuns: sem
vptr, copiáveis, usáveis fora de qualquer cena. E o \texttt{Grupo}, que guarda
\texttt{std::vector<Desenhavel>}, é o padrão \emph{Composite} construído sobre
valores em vez de ponteiros (\cref{sec:patterns}).

\begin{table}[htbp]
\centering
\caption{\texttt{unique\_ptr<Base>} contra o tipo-valor polimórfico.}
\label{tab:valor-vs-ptr}
\footnotesize
\begin{tabular}{@{}lll@{}}
\toprule
\textbf{Aspecto} & \texttt{unique\_ptr<Base>} & \texttt{Desenhavel} \\
\midrule
copiável                      & \textbf{não} & sim \\
o tipo precisa herdar         & sim & \textbf{não} \\
funciona com \texttt{int}/\texttt{string} & não & se atender ao \gls{concept} \\
pode ser \texttt{nullptr}     & sim & \textbf{não} \\
\texttt{==} e \texttt{<}      & compara \textbf{ponteiro} & pode ser por valor \\
aloca                         & sim & sim \\
\bottomrule
\end{tabular}
\end{table}

A linha ``pode ser \texttt{nullptr}'' merece atenção: \texttt{Desenhavel} não tem
construtor padrão, logo não existe ``desenhável vazio'' --- o mesmo argumento de
\texttt{optional<T>} contra \texttt{T*} da \cref{sec:vocabulary}. E a última
linha é a honestidade necessária: os dois alocam e pagam indireção. O ganho aqui
é de \textbf{projeto}, não de velocidade: a medição pôs os dois na mesma faixa,
cerca de 4× o CRTP.

\begin{porque}[por que o \emph{concept} não é enfeite]
Sem o \texttt{SabeDesenhar}, um tipo inadequado produz o erro \emph{dentro} de
\texttt{Modelo<T>}, com cinco linhas de ``\emph{required from}''. Com ele, o
erro volta para a linha do usuário e diz qual expressão faltou: \texttt{note:
the required expression 't.desenhar(nivel)' is invalid}. E o \emph{concept} vira
documentação executável --- ``o que preciso implementar para caber aqui?'' passa
a ter resposta em quatro linhas de código, e não num comentário que envelhece
(\cref{sec:concepts}).
\end{porque}

\begin{armadilha}[quando NÃO fazer apagamento de tipo]
\textbf{(a)}~Você controla todos os tipos e a lista é curta: \texttt{variant} é
4× mais rápido, não aloca e verifica exaustividade --- o apagamento de tipo não
tem exaustividade nenhuma. \textbf{(b)}~O conjunto de operações é grande ou
instável: cada operação nova entra em \emph{três} lugares (Conceito, Modelo e
\emph{wrapper}), e uma interface de vinte métodos vira sessenta linhas de
repasse mecânico. \textbf{(c)}~Não há heterogeneidade em execução: um
\emph{template} resolve em compilação e some. \textbf{(d)}~\texttt{virtual} já
resolve e é mais legível: são cerca de 40 linhas de \emph{boilerplate} para o
que \texttt{virtual} faz em~3.
\end{armadilha}

Existem bibliotecas que geram esse \emph{boilerplate} --- \texttt{dyno}, de
Louis Dionne; \texttt{Boost.TypeErasure}; \texttt{folly::Poly} --- e todas custam
macro ou muita metaprogramação, porque a linguagem ainda não tem reflexão. O
C++26 deve mudar esse quadro.

% =============================================================================
\subsection{Da unidade de tradução ao binário}
\label{sec:linkage}
% =============================================================================

\subsubsection{O que é uma unidade de tradução}

O C++ não compila ``o projeto''. Ele compila \textbf{uma unidade de tradução por
vez}, em total ignorância das outras. Uma \textbf{unidade de tradução} (TU, de
\emph{translation unit}) é \emph{um \texttt{.cpp} mais tudo o que os
\texttt{\#include} dele trouxeram, recursivamente}. Cada TU vira um
arquivo-objeto com duas listas --- os símbolos que ela \textbf{define} e os que
\textbf{precisa} --- e o \emph{linker} cola tudo no fim
(\cref{fig:tu-pipeline}). Praticamente todo erro confuso de \emph{build} é
consequência direta desse modelo.

\begin{figure}[htbp]
\centering
\begin{tikzpicture}[node distance=6mm and 9mm]
  \node[caixa] (a) {\texttt{a.hpp}};
  \node[caixa, below=3mm of a] (m) {\texttt{main.cpp}};
  \node[caixa, below=3mm of m] (b) {\texttt{b.hpp}};
  \node[caixaforte, right=10mm of m] (pp) {pré-processador\\\scriptsize \texttt{g++ -E}};
  \node[caixa, right=9mm of pp] (tu) {uma TU\\\scriptsize um texto só};
  \node[caixaforte, right=9mm of tu] (cc) {compilador\\\scriptsize \texttt{g++ -c}};
  \node[caixa, below=9mm of cc] (o) {\texttt{main.o}\\\scriptsize símbolos};
  \node[caixaforte, left=9mm of o] (ld) {\emph{linker}\\\scriptsize \texttt{ld}};
  \node[caixaok, left=9mm of ld] (exe) {executável};
  \node[caixafraca, below=6mm of ld] (outros) {\texttt{outra.o}, \texttt{libfmt.a}};
  \draw[aresta] (a) -- (pp); \draw[aresta] (m) -- (pp); \draw[aresta] (b) -- (pp);
  \draw[seta] (pp) -- (tu); \draw[seta] (tu) -- (cc); \draw[seta] (cc) -- (o);
  \draw[seta] (o) -- (ld); \draw[seta] (ld) -- (exe); \draw[setafraca] (outros) -- (ld);
\end{tikzpicture}
\caption{O \emph{pipeline} de \texttt{src/linkage/}. O compilador nunca vê duas
         TUs ao mesmo tempo; quem descobre que faltou ou sobrou uma definição é
         o \emph{linker} --- e por isso os erros dele são de outra natureza.}
\label{fig:tu-pipeline}
\end{figure}

Ver cada etapa é uma linha de \emph{shell}: \texttt{g++ -E} para o
pré-processador, \texttt{g++ -c} para o objeto, \texttt{nm -C arq.o} para os
símbolos desmanglados e \texttt{nm -C --undefined-only arq.o} para os que a TU
\emph{precisa}.

\subsubsection{Quanto custa um \texttt{\#include}}

\begin{lstlisting}[style=bt, language=bash,
    caption={Saída de \texttt{src/linkage/}: linhas depois do pré-processador,
             GCC 13.3 com \texttt{-std=c++20}.}]
arquivo vazio       →        6 linhas
<cstddef>           →      253
<vector>            →   26.258
<algorithm>         →   28.270
<string>            →   31.825
<iostream>          →   43.165
<ranges>            →   44.309
\end{lstlisting}

Cada TU paga esse custo \textbf{inteiro}, de novo. Um projeto com cem
\texttt{.cpp} que incluem \cabecalho{string} compila \texttt{<string>} cem
vezes: é por isso que \emph{build} de C++ é lento.

\begin{regra}[reproduza os números que você cita]
A primeira versão desta tabela, no repositório, era \textbf{chutada}: cinco dos
seis números estavam errados, e o \cabecalho{cstddef} errado por um fator de
quatro. Reproduzir custou quatro linhas de \emph{shell} ---
\texttt{echo "\#include <\$h>" > /tmp/t.cpp} seguido de
\texttt{g++ -std=c++20 -E /tmp/t.cpp | wc -l} para cada cabeçalho. Os valores
mudam com a versão da libstdc++, o que é mais um motivo para medir no seu
\emph{toolchain} em vez de copiar os de alguém.
\end{regra}

As saídas, em ordem de eficácia: \textbf{módulos} (C++20 --- a interface é
compilada uma vez e importada em forma binária); \textbf{declaração adiantada}
(\texttt{class Widget;} em vez de \texttt{\#include}, quando só se usa ponteiro
ou referência); e \textbf{pimpl}, assunto da \cref{sec:linkage-abi}.
\emph{Precompiled headers} e \emph{unity builds} atacam o sintoma.

\begin{table}[htbp]
\centering
\caption{Cada mensagem, a etapa em que nasce e a causa típica.}
\label{tab:erros-etapa}
\footnotesize
\begin{tabular}{@{}lll@{}}
\toprule
\textbf{Mensagem} & \textbf{Etapa} & \textbf{Causa típica} \\
\midrule
\texttt{No such file or directory}     & pré-processador & \texttt{-I} faltando \\
\texttt{'X' was not declared}          & compilador & \texttt{\#include} faltando \\
\texttt{invalid use of incomplete type}& compilador & só a \emph{forward declaration} \\
\texttt{undefined reference to 'f()'}  & \textbf{linker} & declarou, não definiu \\
\texttt{multiple definition of 'f()'}  & \textbf{linker} & função não-\texttt{inline}
                                                           em cabeçalho \\
\texttt{undefined reference to vtable} & \textbf{linker} & a \emph{key function}
                                                           não tem corpo \\
\bottomrule
\end{tabular}
\end{table}

\begin{regra}[reconhecer a etapa resolve metade do problema]
Erro do \textbf{compilador} $\rightarrow$ falta uma \textbf{declaração}. Erro do
\textbf{linker} $\rightarrow$ falta (ou sobra) uma \textbf{definição}. Nenhuma
quantidade de \texttt{\#include} conserta um \texttt{undefined reference}.
\end{regra}

A última linha da tabela é a que ninguém adivinha. O GCC emite a \gls{vtable} de
uma classe na TU que define a \textbf{primeira função virtual
não-\texttt{inline} e não-pura} --- a \emph{key function}. Se você declarou
\texttt{virtual void f();} e nunca a definiu, a vtable não é emitida em lugar
nenhum, e o \emph{linker} reclama de um símbolo que você nunca escreveu.

\subsubsection{A ODR, medida com duas TUs reais}

A \textbf{One Definition Rule} tem três partes, e só a terceira surpreende:
\textbf{(1)}~em \emph{uma} TU, cada entidade pode ser declarada muitas vezes e
definida no máximo uma; \textbf{(2)}~no \emph{programa}, cada função ou variável
não-\texttt{inline} pode ser definida no máximo uma vez; \textbf{(3)}~classes,
funções \texttt{inline} e \emph{templates} \textbf{podem} ser definidos em
várias TUs --- mas as definições precisam ser \textbf{token a token idênticas}.
Se não forem, é \gls{ub}, e o \emph{linker} \textbf{não é obrigado a avisar}.

Este é o único subprojeto do repositório com múltiplas unidades de tradução, e
não por capricho. \texttt{unidade\_a.cpp} e \texttt{unidade\_b.cpp} são
idênticos a menos do nome, e cada um devolve os \emph{seus} endereços para as
mesmas entidades de \texttt{comum.hpp}:

\begin{lstlisting}[style=bt, language=C++,
    caption={O palco do experimento (resumo de
             \texttt{src/linkage/comum.hpp}).}]
inline int funcao_inline() { return 1; }   // UMA definição no programa
static int funcao_static() { return 2; }   // uma cópia POR TU
namespace { int funcao_anonima() { return 3; } }   // idem, forma moderna

inline int contador_inline = 0;   // C++17: UM objeto no programa
static int contador_static = 0;   // uma cópia POR TU

const int  kConstante = 42;       // const em namespace: linkage INTERNA
constexpr  int kConstexpr = 43;   // constexpr implica const, logo interna
inline constexpr int kConstexprInline = 44;        // externa

template <typename T>
T identidade(T v) { return v; }   // instanciado em cada TU, e fundido
\end{lstlisting}

\begin{lstlisting}[style=bt, language=bash,
    caption={Saída de \texttt{src/linkage/}: os endereços vêm do \emph{linker},
             não de simulação.}]
inline int f()         MESMO endereço        → 1 no programa inteiro
static int f()         endereços DIFERENTES  → 1 por TU
namespace{} int f()    endereços DIFERENTES  → 1 por TU
inline int var         MESMO endereço        → 1 no programa inteiro
static int var         endereços DIFERENTES  → 1 por TU
const int k = 42       endereços DIFERENTES  → 1 por TU
inline constexpr int   MESMO endereço        → 1 no programa inteiro
template<T> T f(T)     MESMO endereço        → 1 no programa inteiro
\end{lstlisting}

Quatro regras saem daí de uma vez: \texttt{inline} faz o \emph{linker}
\textbf{fundir} as cópias das TUs; \texttt{static} e o \emph{namespace} anônimo
dão \emph{linkage} interna, e cada TU fica com a sua; \texttt{const} em escopo de
\emph{namespace} é implicitamente \texttt{static} em C++ (em C seria externa!);
e um \emph{template} é instanciado em cada TU, com as cópias fundidas como se
fossem \texttt{inline}.

\begin{armadilha}[o \texttt{static} em cabeçalho: N variáveis, não uma]
As duas TUs incrementam \emph{os dois} contadores dez vezes cada:
\begin{lstlisting}[style=bt, language=bash, frame=none,
                   xleftmargin=0pt, xrightmargin=0pt]
| variável         | vista de A | vista de B |
| contador_inline  |         20 |         20 |   ← um objeto
| contador_static  |         10 |         10 |   ← DUAS variáveis
\end{lstlisting}
O bug é traiçoeiro por três motivos: linka sem erro nenhum; funciona
perfeitamente enquanto houver \emph{uma} TU incluindo o cabeçalho; e quebra no
dia em que alguém o inclui num segundo \texttt{.cpp} --- quando o sintoma (``o
\emph{cache} não persiste'') não aponta para o cabeçalho. Antes do C++17 a saída
era \texttt{extern} no cabeçalho mais uma definição num \texttt{.cpp}; desde o
C++17, \texttt{inline} resolve no próprio cabeçalho.
\end{armadilha}

\begin{porque}[\texttt{inline} não é sobre \emph{inlining}]
A palavra engana desde 1998. \texttt{inline} significa \emph{``esta entidade pode
ser definida em várias TUs; o \emph{linker} deve escolher uma e descartar as
outras''}. Ela não obriga o compilador a expandir a chamada --- essa decisão é
do otimizador, que a toma sozinho mesmo em funções não marcadas e a ignora em
funções marcadas que sejam grandes demais. Para \textbf{forçar} a expansão
existe \texttt{[[gnu::always\_inline]]}; para \textbf{impedir},
\texttt{[[gnu::noinline]]}. São implicitamente \texttt{inline}: funções
definidas no corpo da classe, funções \texttt{constexpr} e \texttt{consteval}, e
membros \texttt{static constexpr} (C++17).
\end{porque}

\begin{armadilha}[a violação de ODR que linka, roda e corrompe]
\begin{lstlisting}[style=bt, language=C++, frame=none,
                   xleftmargin=0pt, xrightmargin=0pt]
struct Widget {
    int a;
#ifndef NDEBUG
    int contador_de_depuracao;   // só em debug
#endif
};
// a.cpp SEM -DNDEBUG → sizeof(Widget) == 8
// b.cpp COM -DNDEBUG → sizeof(Widget) == 4
\end{lstlisting}
As duas TUs discordam sobre o \emph{layout} da mesma classe. O \emph{linker} não
compara definições de classe --- ele só vê nomes de símbolo. O programa linka,
roda, e escreve fora dos limites do objeto quando cruza a fronteira. Outras
formas da mesma armadilha: \texttt{-D} diferente entre alvos;
\texttt{\#pragma pack} ativo numa TU e não na outra; \texttt{-std} diferente;
linkar \emph{Debug} com \emph{Release} da mesma biblioteca
(\texttt{\_GLIBCXX\_ASSERTIONS} muda o \emph{layout} dos containers); duas
versões de uma dependência \emph{header-only} na mesma árvore. Defesas:
\textbf{flags uniformes} --- é para isso que serve o \emph{build system}
(\cref{sec:meson}) ---, \texttt{-flto} com \texttt{-Wodr},
\texttt{--detect-odr-violations} no \texttt{gold}/\texttt{lld}, o ASan
(\cref{sec:sanitizers}), e \textbf{nunca} \texttt{\#if} em torno de membro de
classe pública.
\end{armadilha}

\subsubsection{Os três tipos de \emph{linkage}}

\emph{Linkage} responde a uma pergunta só: \textbf{este nome pode se referir à
mesma entidade que um nome igual em outro lugar?} \textbf{Externa}: sim,
inclusive em outra TU. \textbf{Interna}: sim, mas só dentro desta TU.
\textbf{Nenhuma}: não; o nome só existe onde foi declarado.

\begin{regra}[não confunda \emph{linkage}, escopo e tempo de vida]
São três eixos independentes: \emph{linkage} diz quem pode se referir ao nome,
\textbf{escopo} diz onde o nome é visível, \textbf{tempo de vida} diz quando o
objeto existe. Uma variável \texttt{static} dentro de uma função tem
\emph{linkage} \textbf{nenhuma}, escopo de \textbf{bloco} e tempo de vida do
\textbf{programa inteiro} --- três respostas diferentes, na mesma declaração.
\end{regra}

\begin{table}[htbp]
\centering
\caption{\emph{Linkage} por forma de declaração.}
\label{tab:linkage-formas}
\footnotesize
\begin{tabular}{@{}ll@{}}
\toprule
\textbf{Declaração em escopo de \emph{namespace}} & \textbf{\emph{Linkage}} \\
\midrule
\texttt{int x;} \quad \texttt{void f()\{\}} \quad \texttt{inline void f()\{\}} & \textbf{externa} \\
\texttt{static int x;} \quad \texttt{static void f()\{\}} & interna \\
\texttt{const int x = 1;}            & interna \quad --- \emph{a surpresa} \\
\texttt{constexpr int x = 1;}        & interna (implica \texttt{const}) \\
\texttt{inline constexpr int x = 1;} & \textbf{externa} \\
\texttt{namespace \{ int x; \}}      & interna \\
\texttt{extern const int x = 1;}     & \textbf{externa} \\
\texttt{typedef} / \texttt{using} / \emph{enum} sem nome & nenhuma \\
\midrule
\textbf{Declaração em escopo de bloco} & \\
\midrule
\texttt{int x;} (local comum)           & nenhuma \\
\texttt{static int x;} (local estático) & nenhuma --- \emph{mas dura o programa todo} \\
\texttt{extern int x;}                  & \textbf{externa} \\
\bottomrule
\end{tabular}
\end{table}

A linha que mais surpreende é \texttt{const int x = 1;}. Em C ela teria
\emph{linkage} externa; em C++, interna. A razão é prática: permite pôr
constantes em cabeçalhos sem \texttt{multiple definition} --- cada TU ganha a
sua cópia e, como são constantes iguais, ninguém percebe. O preço é o que a
medição mostrou: endereços diferentes em cada TU. Se você precisa de \emph{um}
objeto, escreva \texttt{inline constexpr}.

\subsubsection{\texttt{static}: uma palavra, quatro significados}

\begin{table}[htbp]
\centering
\caption{Os quatro \texttt{static}, que não têm nada em comum entre si.}
\label{tab:quatro-static}
\footnotesize
\begin{tabular}{@{}ll@{}}
\toprule
\textbf{Onde aparece} & \textbf{O que significa} \\
\midrule
escopo de \emph{namespace} & \emph{linkage} \textbf{interna} \\
membro de classe           & \textbf{um por classe}, não por objeto (fora do
                             \texttt{sizeof}) \\
método de classe           & \textbf{sem \texttt{this}}, chamável sem objeto \\
local de função            & \textbf{tempo de vida do programa}, \emph{linkage}
                             nenhuma \\
\bottomrule
\end{tabular}
\end{table}

Quatro conceitos ortogonais na mesma palavra --- herança do C, que o C++ não
pôde consertar sem quebrar tudo. Para o primeiro significado, prefira o
\textbf{\emph{namespace} anônimo}: mesmo efeito, sem ambiguidade, e funciona
também para \textbf{tipos}, coisa que \texttt{static} não faz. O compilador dá
ao \emph{namespace} um nome único por TU e insere um \texttt{using namespace}
implícito; os nomes ficam inacessíveis de fora, e o otimizador ganha uma
informação valiosa --- ninguém de fora chama, então ele pode \emph{inlinar} ou
eliminar à vontade.

\begin{armadilha}[\emph{namespace} anônimo nunca em cabeçalho]
Cada TU que incluir ganha a sua cópia: é o mesmo bug do \texttt{static} em
cabeçalho, agora com \textbf{tipos} --- duas classes com o mesmo nome e
identidades diferentes, o que faz \texttt{dynamic\_cast} e \texttt{typeid}
discordarem entre TUs.
\end{armadilha}

\begin{armadilha}[o \texttt{static} local é \emph{thread-safe}, mas não de graça]
Desde o C++11 a \textbf{inicialização} de um \texttt{static} local é
garantidamente \emph{thread-safe}: se duas \emph{threads} chegarem juntas, uma
inicializa e a outra espera. O subprojeto mede --- 8 \emph{threads} chamando ao
mesmo tempo, \textbf{1 inicialização} ---, e é isso que torna o \emph{Meyers
singleton} correto sem \emph{mutex}. O custo é uma \emph{guard variable} e uma
checagem atômica em \emph{cada} chamada (\texttt{\_\_cxa\_guard\_acquire});
depois da primeira vez é uma leitura atômica relaxada, barata mas não gratuita.
E a garantia é sobre a \emph{inicialização}, não sobre o uso: se várias
\emph{threads} mutam o objeto depois, ainda é preciso sincronizar
(\cref{sec:concorrencia}).
\end{armadilha}

O \texttt{static} local resolve, de quebra, o \emph{static initialization order
fiasco}: globais de TUs diferentes têm ordem de inicialização
\textbf{indefinida} entre si, ao passo que um \texttt{static} local é
inicializado na primeira chamada --- quando tudo de que ele depende já existe.

Quanto a \texttt{extern}, ele diz ``isto é uma \textbf{declaração}, não uma
definição''. Era o único jeito de ter uma global compartilhada antes do C++17;
hoje \texttt{inline int nivel\_de\_log = 3;} faz o mesmo em uma linha. Continua
necessário em três casos: \texttt{extern const int k;} (forçar \emph{linkage}
externa numa constante), \texttt{extern "C"} (\emph{linkage} de
\emph{linguagem}, coisa diferente) e \texttt{extern template} (suprimir a
instanciação numa TU para acelerar o \emph{build}).

\subsubsection{\emph{Name mangling} e \texttt{extern "C"}}

Em C, o símbolo de \texttt{int soma(int,int)} é \texttt{soma}. Em C++ é
\texttt{\_Z4somaii} --- porque o C++ precisa distinguir sobrecargas,
\emph{namespaces} e instanciações de \emph{template}, e o formato de arquivo
objeto só tem \emph{strings} para trabalhar. Essa codificação é o \textbf{name
mangling}.

\begin{lstlisting}[style=bt, language=bash,
    caption={Saída de \texttt{src/linkage/}: os símbolos reais do binário,
             copiados de \texttt{nm}.}]
int f(int)          → _ZN13demo_extern_c8exemplos1fEi
int f(double)       → _ZN13demo_extern_c8exemplos1fEd
int f(int, int)     → _ZN13demo_extern_c8exemplos1fEii
extern "C" f_c(int) → f_c                              ← nome CRU

_Z                prefixo de nome manglado
N ... E           nome aninhado (namespaces)
13demo_extern_c   comprimento + nome
8exemplos         comprimento + nome
1f                comprimento + nome da função
ii                os PARÂMETROS: int, int
\end{lstlisting}

O \emph{mangling} codifica \emph{namespace}, nome e \textbf{tipos dos
parâmetros}. Ele \textbf{não} codifica o tipo de \emph{retorno} --- e é
exatamente por isso que sobrecarregar apenas pelo retorno é proibido na
linguagem: os dois símbolos seriam o mesmo. Para ler à mão, \texttt{nm -C} já
desmangla, e \texttt{echo '\_Z4somaii' | c++filt} resolve casos avulsos.

\begin{armadilha}[o esquema de \emph{mangling} não é padronizado]
GCC e Clang usam o Itanium C++ ABI; o MSVC usa outro, completamente diferente. É
por isso que um \texttt{.lib} do MSVC não linka com o MinGW --- e é o motivo
mais profundo de toda API binária estável do mundo ser exposta em C.
\end{armadilha}

\texttt{extern "C"} diz: ``para esta declaração, use a convenção do C''. Não muda
a linguagem em que o código é escrito --- muda como o \textbf{símbolo} é nomeado.
O subprojeto tem uma API em C \emph{de verdade}, compilada pelo \texttt{cc} e
não pelo \texttt{c++} (o Meson escolhe o compilador pela extensão do arquivo), e
o cabeçalho que os dois lados compartilham traz o idioma canônico:

\begin{lstlisting}[style=bt, language=C++,
    caption={O idioma que todo cabeçalho de biblioteca C traz no topo (resumo de
             \texttt{src/linkage/c\_api.h}).}]
#ifdef __cplusplus
extern "C" {
#endif

int         soma_c(int a, int b);
const char* nome_da_linguagem_c(void);

typedef int (*transformador_t)(int);
int aplica_c(transformador_t f, int v);

#ifdef __cplusplus
}  /* extern "C" */
#endif
\end{lstlisting}

Sem ele, o C++ geraria uma chamada para \texttt{\_Z6soma\_cii} e o \emph{linker}
responderia \texttt{undefined reference to 'soma\_c(int, int)'} --- enquanto
\texttt{nm c\_api.o} mostraria, ali, um \texttt{T soma\_c}. É o erro de
integração C/C++ mais comum que existe, e o \texttt{\_\_cplusplus} é o que
permite ao \emph{mesmo} arquivo servir aos dois compiladores.

\begin{armadilha}[o que \texttt{extern "C"} proíbe --- e o que ele não protege]
\textbf{Não permite}: sobrecarga (há um único símbolo \texttt{soma}, não dois),
\emph{templates} com \emph{linkage} C, funções membro não estáticas.
\textbf{Permite}: usar C++ à vontade \emph{dentro} do corpo --- classes,
exceções, STL. E não protege contra a pior das travessias: \textbf{exceções não
atravessam}. Uma exceção que escapa de uma função \texttt{extern "C"} chamada a
partir de C é \gls{ub} --- as \emph{frames} de C não têm informação de
\emph{unwinding}, e os destrutores intermediários não rodam:
\begin{lstlisting}[style=bt, language=C++, frame=none,
                   xleftmargin=0pt, xrightmargin=0pt]
extern "C" int minha_api(int x) noexcept {
    try { return impl(x); }
    catch (const std::exception& e) { guarda_erro(e.what()); return -1; }
    catch (...) { return -2; }
}
\end{lstlisting}
Erros viram \textbf{códigos de retorno} na fronteira --- é o motivo de toda API
de C ser assim (\cref{sec:errors}). E a outra fronteira que vaza: um objeto
alocado com \texttt{new} no C++ \textbf{não} pode ser liberado com
\texttt{free} no C. Exporte sempre um \texttt{minha\_api\_free(ptr)} junto do
que aloca.
\end{armadilha}

Sobre \emph{callbacks}: um \emph{lambda} sem captura converte para ponteiro de
função, mas o resultado tem \emph{linkage} C++ --- e, a rigor, o tipo de um
ponteiro para função com \emph{linkage} C é diferente do de uma função com
\emph{linkage} C++ (GCC e Clang aceitam os dois; o padrão não garante). Um
\emph{lambda} \textbf{com} captura não converte de jeito nenhum, porque não há
onde guardar o estado; o idioma é o parâmetro \texttt{void* user\_data} que
quase toda API de C oferece, com uma função \texttt{extern "C"} de trampolim.

\subsubsection{ABI: o contrato binário}
\label{sec:linkage-abi}

\begin{porque}[por que a ABI é pior de quebrar do que a API]
\textbf{API} é o contrato de \emph{código-fonte}: quais funções existem, com que
assinatura. Quebrá-la é desagradável e \textbf{seguro} --- o compilador reclama,
alguém conserta, recompila. \textbf{\Gls{abi}} é o contrato \emph{binário}.
Quebrá-la sem quebrar a API é o pior dos mundos: recompila, linka, e
\textbf{corrompe memória em tempo de execução} --- sem uma linha de aviso.
\end{porque}

Fazem parte da ABI mais coisas do que se imagina: os nomes manglados; o
\texttt{sizeof} e o \emph{layout} de cada classe exposta; a \textbf{ordem das
entradas na vtable}; a convenção de chamada; quem destrói cada argumento; o
formato das exceções --- e o \emph{layout} de \textbf{todo tipo da biblioteca
padrão que apareça na sua interface}. Esse último é o que pega: se a sua
biblioteca expõe \texttt{void configura(const std::string\&)}, o \emph{layout} de
\classe{std::string} faz parte da \textbf{sua} ABI. Foi exatamente isso o
rompimento do GCC~5 (\texttt{\_GLIBCXX\_USE\_CXX11\_ABI}): \texttt{std::string}
deixou de ser \emph{copy-on-write} e todo o ecossistema teve de recompilar --- o
símbolo passou a carregar \texttt{\_\_cxx11} no nome justamente para que a
incompatibilidade virasse \textbf{erro de \emph{linker}} em vez de corrupção
silenciosa.

\begin{table}[htbp]
\centering
\caption{Mudanças que quebram a ABI sem quebrar a API.}
\label{tab:abi-quebras}
\footnotesize
\begin{tabular}{@{}lcc@{}}
\toprule
\textbf{Mudança} & \textbf{API} & \textbf{ABI} \\
\midrule
acrescentar ou reordenar membro de dado          & ok & ✗ \\
acrescentar ou reordenar método \textbf{virtual} & ok & ✗ \\
tornar virtual um método comum                   & ok & ✗ \\
acrescentar parâmetro com valor padrão           & ok & ✗ \\
mudar um \emph{enum} de \texttt{int} para \texttt{char} & ok & ✗ \\
tornar \texttt{inline} uma função                & ok & ✗ \\
\midrule
acrescentar método \textbf{não} virtual          & ok & ok \\
acrescentar sobrecarga nova                      & ok & ok \\
acrescentar membro \texttt{static}               & ok & ok \\
mudar o \textbf{corpo} de uma função             & ok & ok \\
\bottomrule
\end{tabular}
\end{table}

Duas linhas surpreendem sistematicamente. ``Acrescentar parâmetro com valor
padrão'': o valor padrão é gravado no lado do \textbf{chamador}, de modo que um
cliente antigo, não recompilado, continua chamando um símbolo que não existe
mais. E ``acrescentar método virtual'': a vtable muda de tamanho e os índices
seguintes se deslocam, então um cliente antigo chamando \texttt{->f()} pega o
\emph{slot} errado e executa outra função, com os argumentos errados. É a
\cref{sec:layout-vptr} cobrando a conta.

\subsubsection{Pimpl: o \emph{firewall} de ABI}

O idioma \emph{pointer to implementation} responde às duas pressões de uma vez:
a classe pública tem um único membro --- um ponteiro --- e tudo o mais vive numa
\texttt{struct} que o cabeçalho nunca vê.

\begin{lstlisting}[style=bt, language=C++,
    caption={A metade pública (resumo de
             \texttt{src/linkage/widget.hpp}).}]
class Widget {
    struct Impl;                                  // só DECLARADA aqui
    std::unique_ptr<Impl> pimpl_;

public:
    explicit Widget(std::string nome);

    // Declarados aqui e definidos em widget.cpp: `unique_ptr<Impl>` precisa
    // do tipo COMPLETO para instanciar o seu deleter.
    ~Widget();
    Widget(Widget&&) noexcept;
    Widget& operator=(Widget&&) noexcept;

    std::string descricao() const;
    void        renomear(std::string novo);
};
\end{lstlisting}

\texttt{sizeof(Widget)} é \textbf{8 bytes}, e continuará 8 para sempre,
independentemente do que a \texttt{Impl} ganhe ou perca. Os ganhos são ABI
estável e \emph{build} rápido --- o cabeçalho não inclui nada do que a
implementação usa, o que ataca diretamente a medição do \texttt{\#include}. Os
custos são uma alocação por objeto, uma indireção por acesso e a impossibilidade
de \emph{inlinar} através da fronteira.

\begin{armadilha}[o destrutor do pimpl não pode ser \texttt{= default} no cabeçalho]
\texttt{\textasciitilde Widget()} precisa do tipo \textbf{completo} de
\texttt{Impl} para instanciar o \emph{deleter} do \classe{std::unique\_ptr}, e no
cabeçalho \texttt{Impl} é incompleta. O erro é críptico ---
\texttt{static assertion failed: can't delete an incomplete type} --- e a
solução é fixa: \textbf{declare} no \texttt{.hpp} e escreva \texttt{= default}
no \texttt{.cpp}, depois da definição de \texttt{Impl}. O mesmo vale para o
construtor e o operador de \gls{move}.
\end{armadilha}

\begin{armadilha}[a segunda armadilha do pimpl custou um erro de \emph{linker} real]
Na primeira versão do subprojeto, \texttt{Widget} estava declarada em
\texttt{demo\_abi.hpp} e \texttt{widget.cpp} incluía esse cabeçalho para
alcançá-la. Como as funções \texttt{demo\_*()} daquele arquivo não são
\texttt{inline}, elas passaram a ser definidas em \textbf{duas} TUs:
\texttt{multiple definition of 'demo\_abi::demo\_o\_que\_e()'}. É o erro da
\cref{tab:erros-etapa}, ao vivo, cometido por quem estava escrevendo sobre ele.
A correção é a que todo projeto acaba adotando: um cabeçalho pequeno, com só o
que precisa ser compartilhado entre TUs, e \textbf{nada definido}.
\end{armadilha}

\paragraph{\emph{Inline namespaces}.} Declarar \texttt{inline namespace v2}
dentro de \texttt{namespace biblioteca} faz os nomes de dentro aparecerem como
se estivessem em \texttt{biblioteca} --- o usuário escreve
\texttt{biblioteca::versao()} sem versão nenhuma ---, mas o símbolo emitido
carrega o \texttt{v2} no \emph{mangling}. Um binário antigo, que pedia
\texttt{v1}, continua encontrando o símbolo de \texttt{v1} se ele ainda existir:
duas versões incompatíveis \textbf{coexistem} no mesmo processo. É como a
libstdc++ faz --- \texttt{std::\_\_cxx11::string} é exatamente isto --- e também
é o mecanismo de \texttt{std::literals::chrono\_literals}.

\begin{regra}[decida cedo o que é ABI no seu código]
Todo mundo sabe que \classe{std::unordered\_map} (um nó alocado por elemento) é
lento comparado a uma tabela de endereçamento aberto, que \classe{std::regex} é
ordens de grandeza mais lento que RE2, e que \texttt{std::vector<bool>} é a
especialização que todos lamentam (\cref{sec:stl}). Consertar qualquer um muda o
\emph{layout}, quebra a ABI e derruba toda biblioteca binária do planeta que os
exponha na interface --- por isso nenhum foi consertado. A lição para o seu
código: se você publica uma biblioteca binária, escolha \textbf{antes} da
primeira versão o que é ABI --- pimpl nas classes que vão crescer, tipos da
biblioteca padrão fora das assinaturas públicas, nenhuma função virtual em
classe que os clientes instanciam. Depois de publicado, a escolha é
irreversível: pergunte à \texttt{std}.
\end{regra}

% =============================================================================
\subsection{Síntese}
% =============================================================================

As três seções parecem tratar de assuntos distantes --- \emph{padding},
\texttt{std::function}, \emph{linker} --- e são três vistas do mesmo objeto: o
\textbf{modelo de objeto} que o compilador implementa por baixo da sintaxe. A
prova é que cada uma explica as armadilhas das outras. O vptr da
\cref{sec:layout} é o que torna o apagamento de tipo por herança
\emph{intrusivo}, e é por isso que a \texttt{std::function} real usa vtable
manual: sem vptr no objeto guardado, ele cabe no \emph{buffer}. O \emph{layout}
de classe \emph{é} metade da \gls{abi} da \cref{sec:linkage} --- acrescentar um
membro de dado e acrescentar um método virtual quebram o binário pela mesma
razão. A medição da \cref{sec:erasure-medicao}, em que a localidade dominou o
despacho por um fator de quatro, é a \cref{sec:layout} explicando por que
reordenar membros vale mais que trocar o mecanismo de chamada. E a substituição
do \texttt{operator new} global, ferramenta de medição do
\projeto{type\_erasure}, só é correta por causa das regras de ODR do
\projeto{linkage}.

Há também um padrão nas medições, e ele é o argumento central deste livro:
\textbf{três das principais surpresas do capítulo contrariaram a expectativa
razoável}. O diamante com herança virtual ficou \emph{maior} tendo \emph{menos}
dados. O derivado com um \texttt{int} a mais teve o \emph{mesmo} \texttt{sizeof}
da base. E o 4× do apagamento de tipo não veio da chamada indireta, mas do cache
--- de modo que a ``otimização'' óbvia, eliminar o \texttt{virtual}, não teria
resolvido nada. Nos três casos, quem raciocinasse sem medir chegaria à conclusão
errada com toda a convicção.

\begin{regra}[as três ferramentas deste capítulo]
\texttt{pahole -C Struct ./binario} mostra os buracos de \emph{padding} que o
\texttt{sizeof} esconde. Um \texttt{operator new} global instrumentado ---
quinze linhas --- responde ``isto aloca?'' para qualquer biblioteca, inclusive a
padrão. E \texttt{nm -C} com \texttt{c++filt} transformam
\texttt{undefined reference} de mistério em pergunta respondível. Nenhuma das
três exige mudar o código de produção, e todas substituem uma discussão por um
número.
\end{regra}
